\section{Voltage and multiple operating points carve the design window}
\label{sec:voltage-design}

The circular EESM map and the PSM interval are loss-limited first screens.
At a prescribed \((M_{\mathrm{req}},\omega_e)\), voltage introduces the set
\begin{equation}
 \Theta_U=\left\{\boldsymbol\theta:\exists\vect i,
 \ M=M_{\mathrm{req}},\ U\leq U_{\max}\right\},
 \qquad
 \boxed{\Theta_{\mathrm{feas}}=\Theta_P\cap\Theta_U.}
 \label{eq:loss-voltage-design-intersection}
\end{equation}
For general EESM parameters \(\Theta_U\) is not a half-plane in the raw
coordinates.  It is exactly and reproducibly represented by the value-function
level set
\begin{equation}
 g_U(\boldsymbol\theta):=
 M_{\max}(P_{\mathrm{Cu},\max},U_{\max},\omega_e;
 \boldsymbol\theta)-M_{\mathrm{req}}\geq0.
 \label{eq:voltage-design-levelset}
\end{equation}
The support-eigenvalue solver of \cref{sec:performance} evaluates this
quantity without a three-current grid.  In a one-parameter design sweep, the
boundary is the scalar root \(g_U=0\); in a two-parameter slice it is a contour;
in the full parameter space it is an implicit hypersurface.  This is a complete
inverse statement even when no useful radical expression exists.

The performance set \(\mathcal W(\boldsymbol\theta)\) offers a second reading.
A machine is feasible when its scaled set contains a direction above the
torque requirement and left of the voltage-per-loss bound.  Changing the
machine deforms the entire convex set: its top point measures maximum torque
per loss, its left-right extent measures voltage sensitivity, and the curvature
of its efficient frontier records their tradeoff.  A small descriptor set can
therefore contain
\begin{equation}
 \left(\lambda_{\max}(\bar H),\ lambda_{\min}^{+}(\bar Q),\
 \lambda_{\max}(\bar Q),\ \text{selected support values of }
 \bar H-\mu\bar Q\right),
 \label{eq:machine-fingerprint}
\end{equation}
where \(\lambda_{\min}^{+}\) is the smallest positive voltage-per-loss
eigenvalue.  The zero eigenvalue of the rank-two EESM voltage matrix is kept
separate because it describes a voltage-null direction, not zero voltage cost
for arbitrary torque.

\subsection{A torque-speed envelope is an intersection, not an average}

For operating points \((\omega_k,M_k)\), define
\begin{equation}
 \gls{sym:thetak}=\left\{\boldsymbol\theta:
 M_{\max}(P_{\mathrm{Cu},\max},U_{\max},\omega_k;
 \boldsymbol\theta)\geq M_k\right\}.
\end{equation}
The complete specification is satisfied if and only if
\begin{equation}
 \boxed{\boldsymbol\theta\in\bigcap_{k=1}^{N}\Theta_k.}
 \label{eq:multipoint-design-intersection}
\end{equation}
A low-speed high-torque point usually pushes outward in torque-per-loss
capability.  A base-speed point adds voltage.  Field-weakening and maximum-speed
points cut away excessive fixed flux or unfavorable voltage-per-loss
directions.  Their common region can be much narrower than the region for any
one point.  Optimizing an average violation would incorrectly admit a machine
that fails a mandatory corner point.

\begin{figure}[tb]
\centering
\begin{tikzpicture}
\begin{axis}[
 width=.82\linewidth,height=.54\linewidth,
 xmin=0,xmax=6,ymin=0,ymax=5,
 xlabel={normalized torque-per-loss coordinate},
 ylabel={field-weakening coordinate},ticks=none,axis lines=left,clip=false]
 \addplot[name path=low,red,very thick,domain=0:5.7] {1.15+4.5/(x+1.25)};
 \addplot[name path=base,blue,very thick,domain=.4:5.8] {.58*x+.35};
 \addplot[name path=high,orange,very thick,domain=.3:5.8] {4.55-.33*x};
 \addplot[name path=floor,draw=none,domain=0:6] {0};
 \addplot[draw=green!50!black,fill=green!28]
   coordinates {(2.55,2.91) (3.25,2.24) (4.35,2.89) (3.67,3.34)} \closedcycle;
 \node[red,anchor=west,font=\scriptsize] at (axis cs:4.1,2.05) {low-speed torque};
 \node[blue,anchor=west,font=\scriptsize,rotate=26] at (axis cs:3.7,2.7) {base speed};
 \node[orange!80!black,anchor=west,font=\scriptsize,rotate=-15] at (axis cs:.7,4.45) {maximum speed};
 \node[green!40!black,align=center,font=\small] at (axis cs:3.45,2.75)
  {feasible\\design window};
 \addplot[only marks,mark=*,black] coordinates {(3.35,2.7)};
 \draw[->,thick] (axis cs:3.35,2.7)--(axis cs:2.95,2.2)
  node[below,font=\scriptsize] {active margin};
\end{axis}
\end{tikzpicture}
\caption{Multiple operating points as intersecting design regions in a
two-dimensional normalized slice.  Each boundary is a level set of the same
current-allocation value function.  Only the common green window satisfies the
entire torque-speed specification.}
\label{fig:multipoint-window}
\end{figure}


The regions in \cref{fig:multipoint-window} are schematic normalized slices;
the actual curves are generated by \cref{eq:voltage-design-levelset}.  This is
also how drive-cycle design studies couple machine parameters to flux-weakening
control: operating-point or cycle constraints are evaluated inside a geometry
search, often with reduced-order models followed by finite-element validation
\cite{dang2017,parasiliti2012}.  The geometric contribution here is to replace
an inner current grid by an exact small value-function evaluation and to retain
the resulting feasibility margins.

\paragraph{Margins are more useful than a Boolean.}
Define \(m_k(\boldsymbol\theta)=M_{\max,k}-M_k\).  The robust nominal margin
is \(m_{\min}=\min_k m_k\).  It identifies the active design point, supports
root bracketing in a chosen parameter, and distinguishes a design barely on
the boundary from one with useful manufacturing and thermal reserve.

